% Supplementary material for the IEEE TAI submission. Build: latexmk -pdf tai_supplement.tex
\documentclass[journal,onecolumn,11pt]{IEEEtran}
\usepackage[T1]{fontenc}
\usepackage{amsmath,amssymb,mathtools}
\let\proof\relax\let\endproof\relax
\usepackage{amsthm}
\usepackage{graphicx,booktabs}
% generated by paper/make_figures.py - do not edit
\newcommand{\NCal}{2179}
\newcommand{\NTest}{1566}
\newcommand{\NQwen}{1424}
\newcommand{\NIntern}{1333}
\newcommand{\NAmber}{1257}
\newcommand{\NSfive}{496}
\newcommand{\NAmberUnknown}{311}
\newcommand{\NAmberClaims}{1568}
\newcommand{\NTestHall}{395}
\newcommand{\NTestImg}{497}
\newcommand{\CleanClaim}{0.080}
\newcommand{\CleanImage}{0.100}
\newcommand{\CleanKept}{0.587}
\newcommand{\GdinoClaim}{0.609 [0.558, 0.657]}
\newcommand{\GdinoImage}{0.409 [0.365, 0.450]}
\newcommand{\GdinoRecClaim}{0.087 [0.060, 0.117]}
\newcommand{\GdinoRecImage}{0.083 [0.060, 0.107]}
\newcommand{\OwlvOneClaim}{0.012 [0.003, 0.023]}
\newcommand{\QwenClaim}{0.129 [0.094, 0.165]}
\newcommand{\InternClaim}{0.160 [0.112, 0.208]}
\newcommand{\QwenImage}{0.093 [0.070, 0.119]}
\newcommand{\InternImage}{0.074 [0.053, 0.098]}
\newcommand{\InternKeptHall}{0.042}
\newcommand{\CleanKeptHall}{0.044}
\newcommand{\AmberClaim}{0.036 [0.008, 0.074]}
\newcommand{\AmberImage}{0.007 [0.002, 0.014]}
\newcommand{\SfiveClaim}{0.105 [0.079, 0.132]}
\newcommand{\SfiveImage}{0.142 [0.112, 0.174]}
\newcommand{\NHtwo}{6}
\newcommand{\Jlearned}{0.14}
\newcommand{\LLgain}{5.45}
\newcommand{\CorrCal}{0.017}
\newcommand{\CorrTest}{0.015}
\newcommand{\AucUnary}{0.870}
\newcommand{\AucIsing}{0.871}
\newcommand{\RobEps}{0.910}
\newcommand{\RobB}{2.711}
\newcommand{\RobImgs}{150}
\newcommand{\WidthJ}{0.506}
\newcommand{\WidthJzero}{0.477}
\newcommand{\TanhJ}{0.139}
\newcommand{\ChairLlava}{0.145}
\newcommand{\ChairQwen}{0.091}
\newcommand{\ChairIntern}{0.076}
\newcommand{\TruncQwen}{47}
\newcommand{\HallRateLlava}{0.252}
\newcommand{\HallAboveLlava}{0.078}
\newcommand{\HallQninetyLlava}{-0.99}
\newcommand{\HallRateQwen}{0.211}
\newcommand{\HallAboveQwen}{0.126}
\newcommand{\HallQninetyQwen}{-0.69}
\newcommand{\HallRateIntern}{0.140}
\newcommand{\HallAboveIntern}{0.155}
\newcommand{\HallQninetyIntern}{-0.47}
\newcommand{\TauFull}{-0.879}
\newcommand{\MinKeptSone}{0.21}
\newcommand{\MinKeptSoneCond}{gaussian noise 5}
\newcommand{\MaxAbstSone}{0.48}
\newcommand{\MaxAbstSoneCond}{gaussian noise 5}
\newcommand{\NClaimPreservedSone}{25}
\newcommand{\NImagePreservedSone}{22}
\newcommand{\PtwoSlack}{0.25}
\newcommand{\ShareBgtEps}{1.00}
\newcommand{\GdinoKept}{0.970}
\newcommand{\GdinoRecKept}{0.599}
\newcommand{\InternKept}{0.602}

\renewcommand{\Pr}{\mathbb{P}}
\begin{document}
\title{Supplementary Material: When Do Conformal Hallucination Guarantees Survive Distribution Shift?}
\author{Anonymous Authors}
\maketitle

\section{Proof of Proposition 1 (Root Dominance)}
Root the tree at the target claim $r$. For a child $c$ of $r$, belief propagation gives the message \[\nu_{c\to r}=\tanh^{-1}(\tanh J_{cr}\tanh\eta_{c\to r}),\] where $\eta_{c\to r}$ is the field of $c$ in the subtree below it, and the field of $r$ under perturbation $\delta$ is $\eta_r(\delta)=\theta_r+\delta_r+\sum_c\nu_{c\to r}(\delta)$; we write $\eta_r=\eta_r(0)$.

\emph{Step 1 (one edge).} For $h(\eta)=\tanh^{-1}(\tanh J\tanh\eta)$ with $J\ge0$,
\[
h'(\eta)=\tanh J\,\frac{1-\tanh^2\eta}{1-\tanh^2J\tanh^2\eta}\in[0,\tanh J],
\]
because $0\le\tanh^2J\le1$ puts the fraction in $[0,1]$. This is the Ising case of the message contraction of Ihler, Fisher and Willsky (JMLR 2005, Theorem 8).

\emph{Step 2 (paths).} The perturbation $\delta_i$ of a claim $i\ne r$ enters $\eta_r(\delta)$ only through the messages along the unique path from $i$ to $r$. By the chain rule, $\partial\eta_r/\partial\delta_i$ is the product of the derivatives of Step 1 along that path, so $0\le\partial\eta_r(\delta)/\partial\delta_i\le\gamma_i$, while $\partial\eta_r(\delta)/\partial\delta_r=1$. In particular $\eta_r(\delta)$ is non-decreasing in every $\delta_i$.

\emph{Step 3 (budget).} By Step 2 the supremum over the uncertainty set is attained with $\delta\ge0$. Writing $\eta_r(\delta)=\eta_r(0,\delta_{-r})+\delta_r$ and applying the mean value theorem on the segment from $0$ to $(0,\delta_{-r})$,
\[
\eta_r(\delta)-\eta_r\le\delta_r+\gamma_{\max}(B-\delta_r),
\]
because the perturbations of the other claims sum to at most $B-\delta_r$. The right side equals $\gamma_{\max}B+(1-\gamma_{\max})\delta_r$ and increases in $\delta_r\le a=\min(\varepsilon_r,B)$, which gives the upper bound; the lower bound is attained by setting $\delta_r=a$ and every other perturbation to zero. If $B\le\varepsilon_r$ the two bounds coincide, the supremum is $\eta_r+B$, and by symmetry the infimum is $\eta_r-B$. Since $\eta\mapsto\sigma(2\eta)$ is increasing, the robust interval for $\Pr(h_r=+1)$ is $[\sigma(2(\eta_r-B)),\sigma(2(\eta_r+B))]$. \hfill$\square$

\emph{Graphical-model analysis.} We fitted a logistic model of each claim's label on its OWLv2 score on Cal-A (300 images), and a chain Ising model over each caption's claims in caption order with one shared coupling $J\in[0,2]$ chosen by exact likelihood on Cal-A. Inference is exact belief propagation. For robust intervals, $\varepsilon$ and $B$ were set to the 0.9 quantiles of the changes in claim fields caused by the corruptions on Cal-A ($\varepsilon=\RobEps$, $B=\RobB$, over \RobImgs{} test images). The coupling improved the Cal-A log-likelihood by \LLgain{} nats over $J=0$. These results are descriptive and test no hypothesis.

\begin{table}[h]\centering\small
\caption{Claim classification (positive class: hallucinated, predicted when the posterior is at least 0.5), unary model ($J=0$) against chain Ising model}
\begin{tabular}{@{}llrrrrrrrrr@{}}\toprule
Data & Model & Acc. & Prec. & Rec. & F1 & AUC & TP & FP & TN & FN\\\midrule
Cal-A & Unary ($J=0$) & 0.820 & 0.753 & 0.502 & 0.602 & 0.859 & 128 & 42 & 641 & 127 \\
Cal-A & Chain Ising (learned $J$) & 0.823 & 0.795 & 0.471 & 0.591 & 0.861 & 120 & 31 & 652 & 135 \\
Test & Unary ($J=0$) & 0.842 & 0.784 & 0.514 & 0.621 & 0.870 & 203 & 56 & 1115 & 192 \\
Test & Chain Ising (learned $J$) & 0.836 & 0.797 & 0.468 & 0.590 & 0.871 & 185 & 47 & 1124 & 210 \\
\bottomrule
\end{tabular}
\end{table}

\section{Results for Every Corruption}
\begin{table}[h]\centering\small
\caption{All S1 conditions (target 0.10): risk [95\% interval], pre-registered decision, share of supported claims kept by the claim-level filter and abstention rate of the image-level filter}
\begin{tabular}{@{}llllllll@{}}\toprule
Corruption & Sev. & Claim-level risk & Decision & Image-level risk & Decision & Kept & Abst.\\\midrule
Gaussian blur & 1 & 0.095 [0.069, 0.125] & preserved & 0.100 [0.076, 0.125] & preserved & 0.616 & 0.025 \\
 & 2 & 0.063 [0.042, 0.087] & preserved & 0.080 [0.057, 0.104] & preserved & 0.594 & 0.030 \\
 & 3 & 0.049 [0.030, 0.072] & preserved & 0.056 [0.037, 0.076] & preserved & 0.549 & 0.045 \\
 & 4 & 0.055 [0.033, 0.077] & preserved & 0.053 [0.034, 0.073] & preserved & 0.504 & 0.065 \\
 & 5 & 0.036 [0.018, 0.054] & preserved & 0.041 [0.024, 0.058] & preserved & 0.408 & 0.155 \\
\addlinespace
JPEG & 1 & 0.096 [0.066, 0.127] & preserved & 0.096 [0.070, 0.120] & preserved & 0.591 & 0.029 \\
 & 2 & 0.067 [0.046, 0.090] & preserved & 0.086 [0.062, 0.108] & preserved & 0.572 & 0.040 \\
 & 3 & 0.063 [0.039, 0.087] & preserved & 0.085 [0.062, 0.108] & preserved & 0.563 & 0.043 \\
 & 4 & 0.044 [0.026, 0.066] & preserved & 0.063 [0.043, 0.084] & preserved & 0.526 & 0.061 \\
 & 5 & 0.039 [0.022, 0.058] & preserved & 0.059 [0.040, 0.080] & preserved & 0.460 & 0.110 \\
\addlinespace
Pixelate & 1 & 0.042 [0.025, 0.062] & preserved & 0.066 [0.045, 0.088] & preserved & 0.540 & 0.048 \\
 & 2 & 0.047 [0.028, 0.068] & preserved & 0.051 [0.033, 0.070] & preserved & 0.542 & 0.050 \\
 & 3 & 0.032 [0.016, 0.049] & preserved & 0.039 [0.023, 0.055] & preserved & 0.512 & 0.072 \\
 & 4 & 0.022 [0.009, 0.037] & preserved & 0.039 [0.023, 0.057] & preserved & 0.413 & 0.153 \\
 & 5 & 0.019 [0.007, 0.033] & preserved & 0.026 [0.013, 0.039] & preserved & 0.340 & 0.258 \\
\addlinespace
Gaussian noise & 1 & 0.092 [0.067, 0.119] & preserved & 0.107 [0.082, 0.133] & \emph{inconcl.} & 0.615 & 0.035 \\
 & 2 & 0.084 [0.058, 0.113] & preserved & 0.089 [0.066, 0.113] & preserved & 0.587 & 0.044 \\
 & 3 & 0.067 [0.045, 0.093] & preserved & 0.080 [0.058, 0.105] & preserved & 0.530 & 0.079 \\
 & 4 & 0.056 [0.035, 0.079] & preserved & 0.062 [0.043, 0.083] & preserved & 0.402 & 0.181 \\
 & 5 & 0.027 [0.012, 0.045] & preserved & 0.027 [0.015, 0.042] & preserved & 0.211 & 0.480 \\
\addlinespace
Contrast & 1 & 0.089 [0.063, 0.118] & preserved & 0.114 [0.088, 0.140] & \emph{inconcl.} & 0.603 & 0.024 \\
 & 2 & 0.096 [0.069, 0.128] & preserved & 0.106 [0.082, 0.132] & \emph{inconcl.} & 0.612 & 0.023 \\
 & 3 & 0.077 [0.053, 0.103] & preserved & 0.097 [0.073, 0.122] & preserved & 0.599 & 0.026 \\
 & 4 & 0.061 [0.041, 0.084] & preserved & 0.075 [0.053, 0.098] & preserved & 0.553 & 0.060 \\
 & 5 & 0.043 [0.024, 0.066] & preserved & 0.048 [0.030, 0.066] & preserved & 0.357 & 0.262 \\
\bottomrule
\end{tabular}
\end{table}

\begin{table}[h]\centering\small
\caption{ImageNet-C severity parameters used for S1}
\begin{tabular}{@{}lrrrrr@{}}\toprule
Corruption & 1 & 2 & 3 & 4 & 5\\\midrule
Gaussian blur, $\sigma$ (pixels) & 1 & 2 & 3 & 4 & 6\\
JPEG quality & 25 & 18 & 15 & 10 & 7\\
Pixelate, resize factor & 0.6 & 0.5 & 0.4 & 0.3 & 0.25\\
Gaussian noise, s.d.\ on $[0,1]$ & 0.08 & 0.12 & 0.18 & 0.26 & 0.38\\
Contrast factor & 0.4 & 0.3 & 0.2 & 0.1 & 0.05\\\bottomrule
\end{tabular}
\end{table}

\section{Deviations from the Pre-Registration}
\begin{table}[h]\centering\small
\begin{tabular}{@{}p{0.27\linewidth}p{0.27\linewidth}p{0.38\linewidth}@{}}\toprule
Planned & Actual & Reason and effect\\\midrule
Proposition 1 (global dominance) as the theoretical account of S1 & Retained as Theorem 1(b); Theorem 1, Theorem 2 and Theorem 3 added after the tests & The confirmatory test of H2 failed while all S1 conditions were preserved; the new results explain this. All analyses that use them are labelled exploratory. Pre-registered decisions are unchanged.\\
Graphical-model analysis in the main text & Moved to Appendix B and this supplement & Descriptive only; it tests no hypothesis.\\\bottomrule
\end{tabular}
\end{table}
\end{document}
